% ============================================================
% GRACE — Appendix (included via \appendix% ============================================================
% GRACE — Appendix (included via \appendix% ============================================================
% GRACE — Appendix (included via \appendix% ============================================================
% GRACE — Appendix (included via \appendix\input{grace_appendix} in grace.tex)
% Not a standalone document — no \documentclass or \begin{document}.
% ============================================================

% Reset counters so appendix figures/tables are numbered A1, A2, ...
\renewcommand{\thefigure}{A\arabic{figure}}
\renewcommand{\thetable}{A\arabic{table}}
\renewcommand{\theequation}{A\arabic{equation}}
\setcounter{figure}{0}
\setcounter{table}{0}
\setcounter{equation}{0}

% ── Placeholder helper (requires tcolorbox — loaded in grace.tex preamble) ────
\newcommand{\figplaceholder}[2]{%
  \begin{tcolorbox}[colback=gray!10,colframe=gray!50,
                    width=\linewidth,height=#1,
                    halign=center,valign=center]
    \textit{\small #2}
  \end{tcolorbox}%
}

\noindent
This appendix provides full implementation details~(\S\ref{sec:impl}),
group-detection evaluation against DBSCAN baselines~(\S\ref{sec:detection}),
extended ablation analysis with training curves and per-seed
breakdowns~(\S\ref{sec:ablation_ext}), a cost-map visualization
gallery~(\S\ref{sec:gallery}), slot assignment
analysis~(\S\ref{sec:slots}), and failure-case analysis~(\S\ref{sec:failures}).

%===============================================================================
\section{Implementation Details}
\label{sec:impl}

\subsection{Architecture Specifications}

Table~\ref{tab:arch} gives the full specification of every module in \grace{}.

\begin{table}[h]
\centering
\caption{Complete \grace{} architecture with parameter counts.}
\label{tab:arch}
\small
\begin{tabular}{llr}
\toprule
\textbf{Module} & \textbf{Specification} & \textbf{Params} \\
\midrule
\multicolumn{3}{l}{\textit{Perception backbone (pretrained, Stage A)}} \\
GroupDetector   & GNN, 3 layers, input 21-d, embed 64-d & \tbd \\
SlotAttention   & $K=3$ slots, 64-d, 3 iters            & \tbd \\
\midrule
\multicolumn{3}{l}{\textit{Cost-map synthesis (analytic, zero learnable params)}} \\
CostMapSynthesizer & 9 ch, $32{\times}32$ grid, 6\,m range & 0 \\
\midrule
\multicolumn{3}{l}{\textit{Planning head (trained, Stages B--C)}} \\
CostMapPlanner  & Conv 32→64→128→128, AdaptivePool(4)        & \tbd \\
Robot MLP       & Linear 9→128→64                            & \tbd \\
Fusion          & Linear 320→256, ReLU                       & \tbd \\
GRUCell         & 256→256                                    & \tbd \\
Actor head      & Linear 256→256, Tanh $\times$2             & \tbd \\
Critic head     & Linear 256→256, Tanh $\times$2, Linear 256→1 & \tbd \\
OccupancyHead   & Conv $C$→32→$T$ ($1{\times}1$), no sigmoid & \tbd \\
\midrule
\textbf{Total trainable (Stage C)} & & \tbd \\
\bottomrule
\end{tabular}
\end{table}

\subsection{Full Hyperparameter Table}

\begin{table}[h]
\centering
\caption{PPO and curriculum hyperparameters for all training stages.}
\label{tab:hparams}
\small
\begin{tabular}{llll}
\toprule
\textbf{Hyperparameter} & \textbf{Stage A} & \textbf{Stage B} & \textbf{Stage C} \\
\midrule
Algorithm           & Supervised          & PPO              & PPO \\
Learning rate       & $1\!\times\!10^{-3}$ & $5\!\times\!10^{-4}$ & $5\!\times\!10^{-5}$ \\
LR schedule         & Cosine              & Linear decay     & Linear decay \\
Parallel envs       & —                   & 16               & 16 \\
Rollout steps       & —                   & 30               & 30 \\
PPO epochs          & —                   & 5                & 5 \\
Mini-batches        & —                   & 2                & 2 \\
Clip $\epsilon$     & —                   & 0.2              & 0.2 \\
Discount $\gamma$   & —                   & 0.99             & 0.99 \\
GAE $\lambda$       & —                   & 0.95             & 0.95 \\
Entropy coef        & —                   & 0.00             & 0.00 \\
Value loss coef     & —                   & 0.5              & 0.5 \\
Grad clip norm      & —                   & 0.5              & 0.5 \\
Aux loss $\lambda$  & —                   & 0.0 (off)        & 0.1 \\
Backbone frozen     & —                   & Yes              & No  \\
Total env steps     & —                   & 20M              & 20M \\
\midrule
\multicolumn{4}{l}{\textit{Environment — Stage B}} \\
\# humans           & — & 10  & — \\
Circle radius (m)   & — & 5.0 & — \\
Groups              & — & 2×static-F & — \\
On-path groups      & — & 0   & — \\
Realistic phases    & — & Off & — \\
\midrule
\multicolumn{4}{l}{\textit{Environment — Stage C}} \\
\# humans           & — & — & 20 \\
Circle radius (m)   & — & — & 8.5 \\
Groups              & — & — & 3×mixed \\
On-path groups      & — & — & 2 \\
Realistic phases    & — & — & A–E (all on) \\
\bottomrule
\end{tabular}
\end{table}

\subsection{Reward Shaping Details}

The reward at each non-terminal step is $r_t = r_\text{pot} + r_\text{grp} + r_\text{ind}$, where:
\begin{align}
r_\text{pot} &= 2\,(\text{dist}_{t-1} - \text{dist}_t) \label{eq:pot} \\
r_\text{ind} &= (d_\text{thr} - d_\text{ind}) \cdot 10 \cdot \Delta t
  \quad \text{if } d_\text{ind} < 0.25\,\text{m} \label{eq:ind} \\
r_\text{grp} &= (d_\text{grp,thr} - d_\text{grp}) \cdot 10 \cdot \Delta t
  \quad \text{if } d_\text{grp} < 0.35\,\text{m} \label{eq:grp}
\end{align}
$r_\text{pot}$ is always active ($\approx{+}0.5$/step at $v_\text{pref}=1$\,m/s).
Terminal events override $r_t$: success $+10$, collision $-20$, group intrusion $-5$.

%===============================================================================
\section{Group Detection Evaluation}
\label{sec:detection}

To validate the GroupDetector + SlotAttention backbone \emph{independently of navigation},
we evaluate it on 100 held-out episodes from the realistic benchmark.
Ground-truth group membership is read from \texttt{human.group\_id};
predicted pair scores are $\hat{y}_{ij} = \sum_k \alpha_{ki}\alpha_{kj}$.

\begin{table}[h]
\centering
\caption{\textbf{Group detection performance vs.\ DBSCAN baselines.}
DBSCAN uses position-only features; DBSCAN+V uses position and velocity.
Averaged over 100 episodes. \tbd{} = run \texttt{eval\_group\_detection.py} to fill.}
\label{tab:detection}
\small
\begin{tabular}{lccccc}
\toprule
\textbf{Method} & \textbf{Pair F1}$\uparrow$ & \textbf{AUC}$\uparrow$
  & \textbf{Purity}$\uparrow$ & \textbf{ARI}$\uparrow$ & \textbf{NMI}$\uparrow$ \\
\midrule
DBSCAN ($\epsilon$=0.5) & \tbd & \tbd & \tbd & \tbd & \tbd \\
DBSCAN ($\epsilon$=0.8) & \tbd & \tbd & \tbd & \tbd & \tbd \\
DBSCAN ($\epsilon$=1.0) & \tbd & \tbd & \tbd & \tbd & \tbd \\
DBSCAN ($\epsilon$=1.5) & \tbd & \tbd & \tbd & \tbd & \tbd \\
DBSCAN+V ($\epsilon$=0.8) & \tbd & \tbd & \tbd & \tbd & \tbd \\
\midrule
\best{\grace{} backbone (ours)} & \tbd & \tbd & \tbd & \tbd & \tbd \\
\bottomrule
\end{tabular}
\end{table}

\textit{Command:}
\texttt{python eval\_group\_detection.py -{}-model\_dir trained\_models/gram\_map/stageC
-{}-test\_model 41665.pt -{}-episodes 100 -{}-seed 0 -{}-out results/detection\_grace.csv}

\subsection{Detection Robustness Under Occlusion}
\label{sec:detection_occlusion}

We repeat the evaluation dropping 20\%, 40\%, and 60\% of visible humans at random
each step to simulate partial occlusion.

\begin{table}[h]
\centering
\caption{\textbf{Pair F1 under increasing occlusion.}
\grace{} is expected to degrade more gracefully than DBSCAN because the GNN
uses velocity context in addition to position.}
\label{tab:occlusion}
\small
\begin{tabular}{lccc}
\toprule
\textbf{Method} & \textbf{0\% dropped} & \textbf{40\% dropped} & \textbf{60\% dropped} \\
\midrule
DBSCAN (best $\epsilon$)  & \tbd & \tbd & \tbd \\
\grace{} backbone         & \tbd & \tbd & \tbd \\
\bottomrule
\end{tabular}
\end{table}

%===============================================================================
\section{Extended Ablation Analysis}
\label{sec:ablation_ext}

\subsection{Training Curves for All Ablation Variants}

\begin{figure}[h]
\centering
\figplaceholder{7cm}{%
  \textbf{Fig.~A\ref{fig:training_curves} — placeholder.}\\[4pt]
  Single plot: mean PPO reward vs.\ training updates.\\
  Lines: \grace{} full (black), C1 no-group (red), C3 no-traj (blue),
  C5 uniform-$\alpha$ (green), C4 no-aux seeds 0/1/2 (orange dashed).\\[4pt]
  \textit{Generate:} \texttt{python plot.py} — configure \texttt{model\_dirs}
  list to include all ablation output directories.\\
  Key visual: C4 orange lines should diverge widely, all others converge.%
}
\caption{\textbf{Training reward curves for \grace{} and all ablation variants.}
  C4 (no auxiliary loss) shows high variance across seeds; all other ablations
  converge, though at lower final reward than the full model.}
\label{fig:training_curves}
\end{figure}

\subsection{Per-Seed Breakdown for C4 (No Auxiliary Loss)}

\begin{table}[h]
\centering
\caption{\textbf{C4: per-seed navigation results.}
Three Stage~C training runs from the same Stage~B checkpoint, auxiliary loss disabled.
High variance confirms the instability claim in Section~6 of the main paper.}
\label{tab:c4_seeds}
\small
\begin{tabular}{lcccc}
\toprule
\textbf{Run} & \textbf{SR\%}$\uparrow$ & \textbf{CR\%}$\downarrow$
  & \textbf{GCR\%}$\downarrow$ & \textbf{Converged?} \\
\midrule
Seed 0              & \tbd & \tbd & \tbd & \tbd \\
Seed 1              & \tbd & \tbd & \tbd & \tbd \\
Seed 2              & \tbd & \tbd & \tbd & \tbd \\
\midrule
Mean $\pm$ std      & \tbd & \tbd & \tbd & — \\
\midrule
\grace{} full (ref) & $87\pm\tbd$ & $7\pm\tbd$ & $0\pm0$ & Yes \\
\bottomrule
\end{tabular}
\end{table}

\subsection{$K$-Slot Sweep}

\begin{figure}[h]
\centering
\figplaceholder{5cm}{%
  \textbf{Fig.~A\ref{fig:k_sweep} — placeholder.}\\[4pt]
  Dual y-axis line plot. x-axis: $K \in \{1,2,3,4,5\}$.
  Left y-axis: SR\% (inverted-U, peak at $K=3$).
  Right y-axis: GCR\% (U-shape, minimum at $K=3$).\\[4pt]
  \textit{Data:} \texttt{test.py} on
  \texttt{trained\_models/gram\_map/ablation\_K\{1,2,4,5\}} + full model ($K=3$).%
}
\caption{\textbf{Effect of slot count $K$ on SR and GCR.}
  $K=3$ achieves the best trade-off: higher SR than $K\leq2$ and lower GCR than
  $K=1$, with marginal degradation at $K>3$ due to over-segmentation.}
\label{fig:k_sweep}
\end{figure}

%===============================================================================
\section{Cost-Map Visualization Gallery}
\label{sec:gallery}

Fig.~\ref{fig:gallery} shows four critical frames from \texttt{visualize\_grace.py}.
Each row is one scenario: (left) ground-truth environment with group hulls;
(middle) robot's learned slot assignments; (right) all 9 cost-map channels.

\begin{figure}[h]
\centering
\figplaceholder{14cm}{%
  \textbf{Fig.~A\ref{fig:gallery} — 4-row gallery placeholder.}\\[6pt]
  Row 1 — Static F-formation blocking path.
    Expected: L6 sharp peak over centroid; robot curves around hull.\\[3pt]
  Row 2 — Tight corridor between two dynamic-LF groups.
    Expected: L7 halos nearly touching; planner threads the gap.\\[3pt]
  Row 3 — Mixed crowd (groups + individuals).
    Expected: middle panel shows 3 distinct slot colors for groups, grey
    for individuals.\\[3pt]
  Row 4 — Failure: bilateral flanking.
    Expected: L6+L7 saturate both sides; robot oscillates.\\[6pt]
  \textit{Generate:}
  \texttt{python visualize\_grace.py -{}-model\_dir trained\_models/gram\_map/stageC
    -{}-test\_model 41665.pt -{}-seed \{3,7,11,15\} -{}-max-steps 400}\\
  Extract key frame: \texttt{ffmpeg -i scene.mp4 -vframes 1 -ss \{T\} frame.png}%
}
\caption{\textbf{Cost-map gallery — four representative scenarios.}
  L6 (group cohesion) consistently locates group centroids;
  L7 (repulsion) provides the early-warning margin that routes the planner
  around group hulls without explicit programming.}
\label{fig:gallery}
\end{figure}

%===============================================================================
\section{Slot Assignment Analysis}
\label{sec:slots}

\subsection{Temporal Stability of Slot Assignments}

\begin{figure}[h]
\centering
\figplaceholder{5cm}{%
  \textbf{Fig.~A\ref{fig:slot_stability} — placeholder.}\\[4pt]
  x-axis: timestep (0–200). y-axis: mean assignment entropy
  $H = -\sum_k \alpha_{kn}\log\alpha_{kn}$ averaged over visible humans.\\
  Expected: high entropy ($\approx$0.9) in steps 0–8, stable below 0.3 afterwards.\\[4pt]
  \textit{Generate:} log \texttt{actor\_critic.base.\_last\_alpha} each step
  in \texttt{test.py}; compute entropy per human per step; plot mean $\pm$ std.%
}
\caption{\textbf{Slot assignment entropy over a full episode.}
  SlotAttention converges to stable, low-entropy group assignments within
  5--10 steps and maintains them for the episode, confirming temporal consistency.}
\label{fig:slot_stability}
\end{figure}

\subsection{Learned Assignments vs.\ Ground-Truth Groups}

\begin{figure}[h]
\centering
\figplaceholder{5cm}{%
  \textbf{Fig.~A\ref{fig:slot_gt} — placeholder.}\\[4pt]
  Side-by-side heatmaps at one timestep. Left: $\alpha \in \mathbb{R}^{3\times20}$
  (rows = slots, columns = humans, color = weight). Right: GT binary membership
  matrix (same layout, one bright column-block per group).\\
  Expected: each GT group block aligns with one dominant slot row on the left.\\[4pt]
  \textit{Extract:} \texttt{actor\_critic.base.\_last\_alpha[0].numpy()} in
  \texttt{visualize\_grace.py}; plot with \texttt{plt.imshow}.%
}
\caption{\textbf{Learned slot assignments vs.\ ground-truth group membership.}
  Each group is captured by one dominant slot; individuals receive distributed
  weights, confirming the competitive-softmax behaviour described in Section~3.2.}
\label{fig:slot_gt}
\end{figure}

%===============================================================================
\section{Failure Case Analysis}
\label{sec:failures}

\subsection{Failure Case 1 — Bilateral Group Corridor}

\begin{figure}[h]
\centering
\figplaceholder{5cm}{%
  \textbf{Fig.~A\ref{fig:failure_bilateral} — placeholder.}\\[4pt]
  Three-panel frame from \texttt{visualize\_grace.py} at the collision step.\\
  Scenario: two dynamic-LF groups converge from both sides simultaneously.\\
  Key: L6 (channel 5) shows two peaks bracketing the robot position;
  L7 (channel 6) halos overlap in the robot's forward direction.\\[4pt]
  \textit{Find seed:} sweep \texttt{-{}-seed} 20–40 until a bilateral-flank
  failure is observed in the left panel trajectory.%
}
\caption{\textbf{Failure case 1: bilateral group corridor.}
  Both L6 peaks are active with no navigable corridor between them.
  The robot oscillates and eventually grazes the nearer group hull.}
\label{fig:failure_bilateral}
\end{figure}

\subsection{Failure Case 2 — Dense Crowd Saturation}

\begin{figure}[h]
\centering
\figplaceholder{5cm}{%
  \textbf{Fig.~A\ref{fig:failure_saturation} — placeholder.}\\[4pt]
  Three-panel frame with \texttt{config.sim.human\_num = 28}.\\
  Key: L1 (individual occupancy) near-uniform ($>$0.8); L6 no longer the
  dominant channel — group cohesion signal is buried under individual noise.\\[4pt]
  \textit{Generate:} temporarily set \texttt{human\_num=28} in config,
  run \texttt{test.py}, capture a timeout episode frame.%
}
\caption{\textbf{Failure case 2: cost-map saturation at $N>25$ pedestrians.}
  Individual and trajectory channels saturate, masking the group cohesion layer.
  The planner reverts to individual-occupancy-based steering and times out.}
\label{fig:failure_saturation}
\end{figure}
 in grace.tex)
% Not a standalone document — no \documentclass or \begin{document}.
% ============================================================

% Reset counters so appendix figures/tables are numbered A1, A2, ...
\renewcommand{\thefigure}{A\arabic{figure}}
\renewcommand{\thetable}{A\arabic{table}}
\renewcommand{\theequation}{A\arabic{equation}}
\setcounter{figure}{0}
\setcounter{table}{0}
\setcounter{equation}{0}

% ── Placeholder helper (requires tcolorbox — loaded in grace.tex preamble) ────
\newcommand{\figplaceholder}[2]{%
  \begin{tcolorbox}[colback=gray!10,colframe=gray!50,
                    width=\linewidth,height=#1,
                    halign=center,valign=center]
    \textit{\small #2}
  \end{tcolorbox}%
}

\noindent
This appendix provides full implementation details~(\S\ref{sec:impl}),
group-detection evaluation against DBSCAN baselines~(\S\ref{sec:detection}),
extended ablation analysis with training curves and per-seed
breakdowns~(\S\ref{sec:ablation_ext}), a cost-map visualization
gallery~(\S\ref{sec:gallery}), slot assignment
analysis~(\S\ref{sec:slots}), and failure-case analysis~(\S\ref{sec:failures}).

%===============================================================================
\section{Implementation Details}
\label{sec:impl}

\subsection{Architecture Specifications}

Table~\ref{tab:arch} gives the full specification of every module in \grace{}.

\begin{table}[h]
\centering
\caption{Complete \grace{} architecture with parameter counts.}
\label{tab:arch}
\small
\begin{tabular}{llr}
\toprule
\textbf{Module} & \textbf{Specification} & \textbf{Params} \\
\midrule
\multicolumn{3}{l}{\textit{Perception backbone (pretrained, Stage A)}} \\
GroupDetector   & GNN, 3 layers, input 21-d, embed 64-d & \tbd \\
SlotAttention   & $K=3$ slots, 64-d, 3 iters            & \tbd \\
\midrule
\multicolumn{3}{l}{\textit{Cost-map synthesis (analytic, zero learnable params)}} \\
CostMapSynthesizer & 9 ch, $32{\times}32$ grid, 6\,m range & 0 \\
\midrule
\multicolumn{3}{l}{\textit{Planning head (trained, Stages B--C)}} \\
CostMapPlanner  & Conv 32→64→128→128, AdaptivePool(4)        & \tbd \\
Robot MLP       & Linear 9→128→64                            & \tbd \\
Fusion          & Linear 320→256, ReLU                       & \tbd \\
GRUCell         & 256→256                                    & \tbd \\
Actor head      & Linear 256→256, Tanh $\times$2             & \tbd \\
Critic head     & Linear 256→256, Tanh $\times$2, Linear 256→1 & \tbd \\
OccupancyHead   & Conv $C$→32→$T$ ($1{\times}1$), no sigmoid & \tbd \\
\midrule
\textbf{Total trainable (Stage C)} & & \tbd \\
\bottomrule
\end{tabular}
\end{table}

\subsection{Full Hyperparameter Table}

\begin{table}[h]
\centering
\caption{PPO and curriculum hyperparameters for all training stages.}
\label{tab:hparams}
\small
\begin{tabular}{llll}
\toprule
\textbf{Hyperparameter} & \textbf{Stage A} & \textbf{Stage B} & \textbf{Stage C} \\
\midrule
Algorithm           & Supervised          & PPO              & PPO \\
Learning rate       & $1\!\times\!10^{-3}$ & $5\!\times\!10^{-4}$ & $5\!\times\!10^{-5}$ \\
LR schedule         & Cosine              & Linear decay     & Linear decay \\
Parallel envs       & —                   & 16               & 16 \\
Rollout steps       & —                   & 30               & 30 \\
PPO epochs          & —                   & 5                & 5 \\
Mini-batches        & —                   & 2                & 2 \\
Clip $\epsilon$     & —                   & 0.2              & 0.2 \\
Discount $\gamma$   & —                   & 0.99             & 0.99 \\
GAE $\lambda$       & —                   & 0.95             & 0.95 \\
Entropy coef        & —                   & 0.00             & 0.00 \\
Value loss coef     & —                   & 0.5              & 0.5 \\
Grad clip norm      & —                   & 0.5              & 0.5 \\
Aux loss $\lambda$  & —                   & 0.0 (off)        & 0.1 \\
Backbone frozen     & —                   & Yes              & No  \\
Total env steps     & —                   & 20M              & 20M \\
\midrule
\multicolumn{4}{l}{\textit{Environment — Stage B}} \\
\# humans           & — & 10  & — \\
Circle radius (m)   & — & 5.0 & — \\
Groups              & — & 2×static-F & — \\
On-path groups      & — & 0   & — \\
Realistic phases    & — & Off & — \\
\midrule
\multicolumn{4}{l}{\textit{Environment — Stage C}} \\
\# humans           & — & — & 20 \\
Circle radius (m)   & — & — & 8.5 \\
Groups              & — & — & 3×mixed \\
On-path groups      & — & — & 2 \\
Realistic phases    & — & — & A–E (all on) \\
\bottomrule
\end{tabular}
\end{table}

\subsection{Reward Shaping Details}

The reward at each non-terminal step is $r_t = r_\text{pot} + r_\text{grp} + r_\text{ind}$, where:
\begin{align}
r_\text{pot} &= 2\,(\text{dist}_{t-1} - \text{dist}_t) \label{eq:pot} \\
r_\text{ind} &= (d_\text{thr} - d_\text{ind}) \cdot 10 \cdot \Delta t
  \quad \text{if } d_\text{ind} < 0.25\,\text{m} \label{eq:ind} \\
r_\text{grp} &= (d_\text{grp,thr} - d_\text{grp}) \cdot 10 \cdot \Delta t
  \quad \text{if } d_\text{grp} < 0.35\,\text{m} \label{eq:grp}
\end{align}
$r_\text{pot}$ is always active ($\approx{+}0.5$/step at $v_\text{pref}=1$\,m/s).
Terminal events override $r_t$: success $+10$, collision $-20$, group intrusion $-5$.

%===============================================================================
\section{Group Detection Evaluation}
\label{sec:detection}

To validate the GroupDetector + SlotAttention backbone \emph{independently of navigation},
we evaluate it on 100 held-out episodes from the realistic benchmark.
Ground-truth group membership is read from \texttt{human.group\_id};
predicted pair scores are $\hat{y}_{ij} = \sum_k \alpha_{ki}\alpha_{kj}$.

\begin{table}[h]
\centering
\caption{\textbf{Group detection performance vs.\ DBSCAN baselines.}
DBSCAN uses position-only features; DBSCAN+V uses position and velocity.
Averaged over 100 episodes. \tbd{} = run \texttt{eval\_group\_detection.py} to fill.}
\label{tab:detection}
\small
\begin{tabular}{lccccc}
\toprule
\textbf{Method} & \textbf{Pair F1}$\uparrow$ & \textbf{AUC}$\uparrow$
  & \textbf{Purity}$\uparrow$ & \textbf{ARI}$\uparrow$ & \textbf{NMI}$\uparrow$ \\
\midrule
DBSCAN ($\epsilon$=0.5) & \tbd & \tbd & \tbd & \tbd & \tbd \\
DBSCAN ($\epsilon$=0.8) & \tbd & \tbd & \tbd & \tbd & \tbd \\
DBSCAN ($\epsilon$=1.0) & \tbd & \tbd & \tbd & \tbd & \tbd \\
DBSCAN ($\epsilon$=1.5) & \tbd & \tbd & \tbd & \tbd & \tbd \\
DBSCAN+V ($\epsilon$=0.8) & \tbd & \tbd & \tbd & \tbd & \tbd \\
\midrule
\best{\grace{} backbone (ours)} & \tbd & \tbd & \tbd & \tbd & \tbd \\
\bottomrule
\end{tabular}
\end{table}

\textit{Command:}
\texttt{python eval\_group\_detection.py -{}-model\_dir trained\_models/gram\_map/stageC
-{}-test\_model 41665.pt -{}-episodes 100 -{}-seed 0 -{}-out results/detection\_grace.csv}

\subsection{Detection Robustness Under Occlusion}
\label{sec:detection_occlusion}

We repeat the evaluation dropping 20\%, 40\%, and 60\% of visible humans at random
each step to simulate partial occlusion.

\begin{table}[h]
\centering
\caption{\textbf{Pair F1 under increasing occlusion.}
\grace{} is expected to degrade more gracefully than DBSCAN because the GNN
uses velocity context in addition to position.}
\label{tab:occlusion}
\small
\begin{tabular}{lccc}
\toprule
\textbf{Method} & \textbf{0\% dropped} & \textbf{40\% dropped} & \textbf{60\% dropped} \\
\midrule
DBSCAN (best $\epsilon$)  & \tbd & \tbd & \tbd \\
\grace{} backbone         & \tbd & \tbd & \tbd \\
\bottomrule
\end{tabular}
\end{table}

%===============================================================================
\section{Extended Ablation Analysis}
\label{sec:ablation_ext}

\subsection{Training Curves for All Ablation Variants}

\begin{figure}[h]
\centering
\figplaceholder{7cm}{%
  \textbf{Fig.~A\ref{fig:training_curves} — placeholder.}\\[4pt]
  Single plot: mean PPO reward vs.\ training updates.\\
  Lines: \grace{} full (black), C1 no-group (red), C3 no-traj (blue),
  C5 uniform-$\alpha$ (green), C4 no-aux seeds 0/1/2 (orange dashed).\\[4pt]
  \textit{Generate:} \texttt{python plot.py} — configure \texttt{model\_dirs}
  list to include all ablation output directories.\\
  Key visual: C4 orange lines should diverge widely, all others converge.%
}
\caption{\textbf{Training reward curves for \grace{} and all ablation variants.}
  C4 (no auxiliary loss) shows high variance across seeds; all other ablations
  converge, though at lower final reward than the full model.}
\label{fig:training_curves}
\end{figure}

\subsection{Per-Seed Breakdown for C4 (No Auxiliary Loss)}

\begin{table}[h]
\centering
\caption{\textbf{C4: per-seed navigation results.}
Three Stage~C training runs from the same Stage~B checkpoint, auxiliary loss disabled.
High variance confirms the instability claim in Section~6 of the main paper.}
\label{tab:c4_seeds}
\small
\begin{tabular}{lcccc}
\toprule
\textbf{Run} & \textbf{SR\%}$\uparrow$ & \textbf{CR\%}$\downarrow$
  & \textbf{GCR\%}$\downarrow$ & \textbf{Converged?} \\
\midrule
Seed 0              & \tbd & \tbd & \tbd & \tbd \\
Seed 1              & \tbd & \tbd & \tbd & \tbd \\
Seed 2              & \tbd & \tbd & \tbd & \tbd \\
\midrule
Mean $\pm$ std      & \tbd & \tbd & \tbd & — \\
\midrule
\grace{} full (ref) & $87\pm\tbd$ & $7\pm\tbd$ & $0\pm0$ & Yes \\
\bottomrule
\end{tabular}
\end{table}

\subsection{$K$-Slot Sweep}

\begin{figure}[h]
\centering
\figplaceholder{5cm}{%
  \textbf{Fig.~A\ref{fig:k_sweep} — placeholder.}\\[4pt]
  Dual y-axis line plot. x-axis: $K \in \{1,2,3,4,5\}$.
  Left y-axis: SR\% (inverted-U, peak at $K=3$).
  Right y-axis: GCR\% (U-shape, minimum at $K=3$).\\[4pt]
  \textit{Data:} \texttt{test.py} on
  \texttt{trained\_models/gram\_map/ablation\_K\{1,2,4,5\}} + full model ($K=3$).%
}
\caption{\textbf{Effect of slot count $K$ on SR and GCR.}
  $K=3$ achieves the best trade-off: higher SR than $K\leq2$ and lower GCR than
  $K=1$, with marginal degradation at $K>3$ due to over-segmentation.}
\label{fig:k_sweep}
\end{figure}

%===============================================================================
\section{Cost-Map Visualization Gallery}
\label{sec:gallery}

Fig.~\ref{fig:gallery} shows four critical frames from \texttt{visualize\_grace.py}.
Each row is one scenario: (left) ground-truth environment with group hulls;
(middle) robot's learned slot assignments; (right) all 9 cost-map channels.

\begin{figure}[h]
\centering
\figplaceholder{14cm}{%
  \textbf{Fig.~A\ref{fig:gallery} — 4-row gallery placeholder.}\\[6pt]
  Row 1 — Static F-formation blocking path.
    Expected: L6 sharp peak over centroid; robot curves around hull.\\[3pt]
  Row 2 — Tight corridor between two dynamic-LF groups.
    Expected: L7 halos nearly touching; planner threads the gap.\\[3pt]
  Row 3 — Mixed crowd (groups + individuals).
    Expected: middle panel shows 3 distinct slot colors for groups, grey
    for individuals.\\[3pt]
  Row 4 — Failure: bilateral flanking.
    Expected: L6+L7 saturate both sides; robot oscillates.\\[6pt]
  \textit{Generate:}
  \texttt{python visualize\_grace.py -{}-model\_dir trained\_models/gram\_map/stageC
    -{}-test\_model 41665.pt -{}-seed \{3,7,11,15\} -{}-max-steps 400}\\
  Extract key frame: \texttt{ffmpeg -i scene.mp4 -vframes 1 -ss \{T\} frame.png}%
}
\caption{\textbf{Cost-map gallery — four representative scenarios.}
  L6 (group cohesion) consistently locates group centroids;
  L7 (repulsion) provides the early-warning margin that routes the planner
  around group hulls without explicit programming.}
\label{fig:gallery}
\end{figure}

%===============================================================================
\section{Slot Assignment Analysis}
\label{sec:slots}

\subsection{Temporal Stability of Slot Assignments}

\begin{figure}[h]
\centering
\figplaceholder{5cm}{%
  \textbf{Fig.~A\ref{fig:slot_stability} — placeholder.}\\[4pt]
  x-axis: timestep (0–200). y-axis: mean assignment entropy
  $H = -\sum_k \alpha_{kn}\log\alpha_{kn}$ averaged over visible humans.\\
  Expected: high entropy ($\approx$0.9) in steps 0–8, stable below 0.3 afterwards.\\[4pt]
  \textit{Generate:} log \texttt{actor\_critic.base.\_last\_alpha} each step
  in \texttt{test.py}; compute entropy per human per step; plot mean $\pm$ std.%
}
\caption{\textbf{Slot assignment entropy over a full episode.}
  SlotAttention converges to stable, low-entropy group assignments within
  5--10 steps and maintains them for the episode, confirming temporal consistency.}
\label{fig:slot_stability}
\end{figure}

\subsection{Learned Assignments vs.\ Ground-Truth Groups}

\begin{figure}[h]
\centering
\figplaceholder{5cm}{%
  \textbf{Fig.~A\ref{fig:slot_gt} — placeholder.}\\[4pt]
  Side-by-side heatmaps at one timestep. Left: $\alpha \in \mathbb{R}^{3\times20}$
  (rows = slots, columns = humans, color = weight). Right: GT binary membership
  matrix (same layout, one bright column-block per group).\\
  Expected: each GT group block aligns with one dominant slot row on the left.\\[4pt]
  \textit{Extract:} \texttt{actor\_critic.base.\_last\_alpha[0].numpy()} in
  \texttt{visualize\_grace.py}; plot with \texttt{plt.imshow}.%
}
\caption{\textbf{Learned slot assignments vs.\ ground-truth group membership.}
  Each group is captured by one dominant slot; individuals receive distributed
  weights, confirming the competitive-softmax behaviour described in Section~3.2.}
\label{fig:slot_gt}
\end{figure}

%===============================================================================
\section{Failure Case Analysis}
\label{sec:failures}

\subsection{Failure Case 1 — Bilateral Group Corridor}

\begin{figure}[h]
\centering
\figplaceholder{5cm}{%
  \textbf{Fig.~A\ref{fig:failure_bilateral} — placeholder.}\\[4pt]
  Three-panel frame from \texttt{visualize\_grace.py} at the collision step.\\
  Scenario: two dynamic-LF groups converge from both sides simultaneously.\\
  Key: L6 (channel 5) shows two peaks bracketing the robot position;
  L7 (channel 6) halos overlap in the robot's forward direction.\\[4pt]
  \textit{Find seed:} sweep \texttt{-{}-seed} 20–40 until a bilateral-flank
  failure is observed in the left panel trajectory.%
}
\caption{\textbf{Failure case 1: bilateral group corridor.}
  Both L6 peaks are active with no navigable corridor between them.
  The robot oscillates and eventually grazes the nearer group hull.}
\label{fig:failure_bilateral}
\end{figure}

\subsection{Failure Case 2 — Dense Crowd Saturation}

\begin{figure}[h]
\centering
\figplaceholder{5cm}{%
  \textbf{Fig.~A\ref{fig:failure_saturation} — placeholder.}\\[4pt]
  Three-panel frame with \texttt{config.sim.human\_num = 28}.\\
  Key: L1 (individual occupancy) near-uniform ($>$0.8); L6 no longer the
  dominant channel — group cohesion signal is buried under individual noise.\\[4pt]
  \textit{Generate:} temporarily set \texttt{human\_num=28} in config,
  run \texttt{test.py}, capture a timeout episode frame.%
}
\caption{\textbf{Failure case 2: cost-map saturation at $N>25$ pedestrians.}
  Individual and trajectory channels saturate, masking the group cohesion layer.
  The planner reverts to individual-occupancy-based steering and times out.}
\label{fig:failure_saturation}
\end{figure}
 in grace.tex)
% Not a standalone document — no \documentclass or \begin{document}.
% ============================================================

% Reset counters so appendix figures/tables are numbered A1, A2, ...
\renewcommand{\thefigure}{A\arabic{figure}}
\renewcommand{\thetable}{A\arabic{table}}
\renewcommand{\theequation}{A\arabic{equation}}
\setcounter{figure}{0}
\setcounter{table}{0}
\setcounter{equation}{0}

% ── Placeholder helper (requires tcolorbox — loaded in grace.tex preamble) ────
\newcommand{\figplaceholder}[2]{%
  \begin{tcolorbox}[colback=gray!10,colframe=gray!50,
                    width=\linewidth,height=#1,
                    halign=center,valign=center]
    \textit{\small #2}
  \end{tcolorbox}%
}

\noindent
This appendix provides full implementation details~(\S\ref{sec:impl}),
group-detection evaluation against DBSCAN baselines~(\S\ref{sec:detection}),
extended ablation analysis with training curves and per-seed
breakdowns~(\S\ref{sec:ablation_ext}), a cost-map visualization
gallery~(\S\ref{sec:gallery}), slot assignment
analysis~(\S\ref{sec:slots}), and failure-case analysis~(\S\ref{sec:failures}).

%===============================================================================
\section{Implementation Details}
\label{sec:impl}

\subsection{Architecture Specifications}

Table~\ref{tab:arch} gives the full specification of every module in \grace{}.

\begin{table}[h]
\centering
\caption{Complete \grace{} architecture with parameter counts.}
\label{tab:arch}
\small
\begin{tabular}{llr}
\toprule
\textbf{Module} & \textbf{Specification} & \textbf{Params} \\
\midrule
\multicolumn{3}{l}{\textit{Perception backbone (pretrained, Stage A)}} \\
GroupDetector   & GNN, 3 layers, input 21-d, embed 64-d & \tbd \\
SlotAttention   & $K=3$ slots, 64-d, 3 iters            & \tbd \\
\midrule
\multicolumn{3}{l}{\textit{Cost-map synthesis (analytic, zero learnable params)}} \\
CostMapSynthesizer & 9 ch, $32{\times}32$ grid, 6\,m range & 0 \\
\midrule
\multicolumn{3}{l}{\textit{Planning head (trained, Stages B--C)}} \\
CostMapPlanner  & Conv 32→64→128→128, AdaptivePool(4)        & \tbd \\
Robot MLP       & Linear 9→128→64                            & \tbd \\
Fusion          & Linear 320→256, ReLU                       & \tbd \\
GRUCell         & 256→256                                    & \tbd \\
Actor head      & Linear 256→256, Tanh $\times$2             & \tbd \\
Critic head     & Linear 256→256, Tanh $\times$2, Linear 256→1 & \tbd \\
OccupancyHead   & Conv $C$→32→$T$ ($1{\times}1$), no sigmoid & \tbd \\
\midrule
\textbf{Total trainable (Stage C)} & & \tbd \\
\bottomrule
\end{tabular}
\end{table}

\subsection{Full Hyperparameter Table}

\begin{table}[h]
\centering
\caption{PPO and curriculum hyperparameters for all training stages.}
\label{tab:hparams}
\small
\begin{tabular}{llll}
\toprule
\textbf{Hyperparameter} & \textbf{Stage A} & \textbf{Stage B} & \textbf{Stage C} \\
\midrule
Algorithm           & Supervised          & PPO              & PPO \\
Learning rate       & $1\!\times\!10^{-3}$ & $5\!\times\!10^{-4}$ & $5\!\times\!10^{-5}$ \\
LR schedule         & Cosine              & Linear decay     & Linear decay \\
Parallel envs       & —                   & 16               & 16 \\
Rollout steps       & —                   & 30               & 30 \\
PPO epochs          & —                   & 5                & 5 \\
Mini-batches        & —                   & 2                & 2 \\
Clip $\epsilon$     & —                   & 0.2              & 0.2 \\
Discount $\gamma$   & —                   & 0.99             & 0.99 \\
GAE $\lambda$       & —                   & 0.95             & 0.95 \\
Entropy coef        & —                   & 0.00             & 0.00 \\
Value loss coef     & —                   & 0.5              & 0.5 \\
Grad clip norm      & —                   & 0.5              & 0.5 \\
Aux loss $\lambda$  & —                   & 0.0 (off)        & 0.1 \\
Backbone frozen     & —                   & Yes              & No  \\
Total env steps     & —                   & 20M              & 20M \\
\midrule
\multicolumn{4}{l}{\textit{Environment — Stage B}} \\
\# humans           & — & 10  & — \\
Circle radius (m)   & — & 5.0 & — \\
Groups              & — & 2×static-F & — \\
On-path groups      & — & 0   & — \\
Realistic phases    & — & Off & — \\
\midrule
\multicolumn{4}{l}{\textit{Environment — Stage C}} \\
\# humans           & — & — & 20 \\
Circle radius (m)   & — & — & 8.5 \\
Groups              & — & — & 3×mixed \\
On-path groups      & — & — & 2 \\
Realistic phases    & — & — & A–E (all on) \\
\bottomrule
\end{tabular}
\end{table}

\subsection{Reward Shaping Details}

The reward at each non-terminal step is $r_t = r_\text{pot} + r_\text{grp} + r_\text{ind}$, where:
\begin{align}
r_\text{pot} &= 2\,(\text{dist}_{t-1} - \text{dist}_t) \label{eq:pot} \\
r_\text{ind} &= (d_\text{thr} - d_\text{ind}) \cdot 10 \cdot \Delta t
  \quad \text{if } d_\text{ind} < 0.25\,\text{m} \label{eq:ind} \\
r_\text{grp} &= (d_\text{grp,thr} - d_\text{grp}) \cdot 10 \cdot \Delta t
  \quad \text{if } d_\text{grp} < 0.35\,\text{m} \label{eq:grp}
\end{align}
$r_\text{pot}$ is always active ($\approx{+}0.5$/step at $v_\text{pref}=1$\,m/s).
Terminal events override $r_t$: success $+10$, collision $-20$, group intrusion $-5$.

%===============================================================================
\section{Group Detection Evaluation}
\label{sec:detection}

To validate the GroupDetector + SlotAttention backbone \emph{independently of navigation},
we evaluate it on 100 held-out episodes from the realistic benchmark.
Ground-truth group membership is read from \texttt{human.group\_id};
predicted pair scores are $\hat{y}_{ij} = \sum_k \alpha_{ki}\alpha_{kj}$.

\begin{table}[h]
\centering
\caption{\textbf{Group detection performance vs.\ DBSCAN baselines.}
DBSCAN uses position-only features; DBSCAN+V uses position and velocity.
Averaged over 100 episodes. \tbd{} = run \texttt{eval\_group\_detection.py} to fill.}
\label{tab:detection}
\small
\begin{tabular}{lccccc}
\toprule
\textbf{Method} & \textbf{Pair F1}$\uparrow$ & \textbf{AUC}$\uparrow$
  & \textbf{Purity}$\uparrow$ & \textbf{ARI}$\uparrow$ & \textbf{NMI}$\uparrow$ \\
\midrule
DBSCAN ($\epsilon$=0.5) & \tbd & \tbd & \tbd & \tbd & \tbd \\
DBSCAN ($\epsilon$=0.8) & \tbd & \tbd & \tbd & \tbd & \tbd \\
DBSCAN ($\epsilon$=1.0) & \tbd & \tbd & \tbd & \tbd & \tbd \\
DBSCAN ($\epsilon$=1.5) & \tbd & \tbd & \tbd & \tbd & \tbd \\
DBSCAN+V ($\epsilon$=0.8) & \tbd & \tbd & \tbd & \tbd & \tbd \\
\midrule
\best{\grace{} backbone (ours)} & \tbd & \tbd & \tbd & \tbd & \tbd \\
\bottomrule
\end{tabular}
\end{table}

\textit{Command:}
\texttt{python eval\_group\_detection.py -{}-model\_dir trained\_models/gram\_map/stageC
-{}-test\_model 41665.pt -{}-episodes 100 -{}-seed 0 -{}-out results/detection\_grace.csv}

\subsection{Detection Robustness Under Occlusion}
\label{sec:detection_occlusion}

We repeat the evaluation dropping 20\%, 40\%, and 60\% of visible humans at random
each step to simulate partial occlusion.

\begin{table}[h]
\centering
\caption{\textbf{Pair F1 under increasing occlusion.}
\grace{} is expected to degrade more gracefully than DBSCAN because the GNN
uses velocity context in addition to position.}
\label{tab:occlusion}
\small
\begin{tabular}{lccc}
\toprule
\textbf{Method} & \textbf{0\% dropped} & \textbf{40\% dropped} & \textbf{60\% dropped} \\
\midrule
DBSCAN (best $\epsilon$)  & \tbd & \tbd & \tbd \\
\grace{} backbone         & \tbd & \tbd & \tbd \\
\bottomrule
\end{tabular}
\end{table}

%===============================================================================
\section{Extended Ablation Analysis}
\label{sec:ablation_ext}

\subsection{Training Curves for All Ablation Variants}

\begin{figure}[h]
\centering
\figplaceholder{7cm}{%
  \textbf{Fig.~A\ref{fig:training_curves} — placeholder.}\\[4pt]
  Single plot: mean PPO reward vs.\ training updates.\\
  Lines: \grace{} full (black), C1 no-group (red), C3 no-traj (blue),
  C5 uniform-$\alpha$ (green), C4 no-aux seeds 0/1/2 (orange dashed).\\[4pt]
  \textit{Generate:} \texttt{python plot.py} — configure \texttt{model\_dirs}
  list to include all ablation output directories.\\
  Key visual: C4 orange lines should diverge widely, all others converge.%
}
\caption{\textbf{Training reward curves for \grace{} and all ablation variants.}
  C4 (no auxiliary loss) shows high variance across seeds; all other ablations
  converge, though at lower final reward than the full model.}
\label{fig:training_curves}
\end{figure}

\subsection{Per-Seed Breakdown for C4 (No Auxiliary Loss)}

\begin{table}[h]
\centering
\caption{\textbf{C4: per-seed navigation results.}
Three Stage~C training runs from the same Stage~B checkpoint, auxiliary loss disabled.
High variance confirms the instability claim in Section~6 of the main paper.}
\label{tab:c4_seeds}
\small
\begin{tabular}{lcccc}
\toprule
\textbf{Run} & \textbf{SR\%}$\uparrow$ & \textbf{CR\%}$\downarrow$
  & \textbf{GCR\%}$\downarrow$ & \textbf{Converged?} \\
\midrule
Seed 0              & \tbd & \tbd & \tbd & \tbd \\
Seed 1              & \tbd & \tbd & \tbd & \tbd \\
Seed 2              & \tbd & \tbd & \tbd & \tbd \\
\midrule
Mean $\pm$ std      & \tbd & \tbd & \tbd & — \\
\midrule
\grace{} full (ref) & $87\pm\tbd$ & $7\pm\tbd$ & $0\pm0$ & Yes \\
\bottomrule
\end{tabular}
\end{table}

\subsection{$K$-Slot Sweep}

\begin{figure}[h]
\centering
\figplaceholder{5cm}{%
  \textbf{Fig.~A\ref{fig:k_sweep} — placeholder.}\\[4pt]
  Dual y-axis line plot. x-axis: $K \in \{1,2,3,4,5\}$.
  Left y-axis: SR\% (inverted-U, peak at $K=3$).
  Right y-axis: GCR\% (U-shape, minimum at $K=3$).\\[4pt]
  \textit{Data:} \texttt{test.py} on
  \texttt{trained\_models/gram\_map/ablation\_K\{1,2,4,5\}} + full model ($K=3$).%
}
\caption{\textbf{Effect of slot count $K$ on SR and GCR.}
  $K=3$ achieves the best trade-off: higher SR than $K\leq2$ and lower GCR than
  $K=1$, with marginal degradation at $K>3$ due to over-segmentation.}
\label{fig:k_sweep}
\end{figure}

%===============================================================================
\section{Cost-Map Visualization Gallery}
\label{sec:gallery}

Fig.~\ref{fig:gallery} shows four critical frames from \texttt{visualize\_grace.py}.
Each row is one scenario: (left) ground-truth environment with group hulls;
(middle) robot's learned slot assignments; (right) all 9 cost-map channels.

\begin{figure}[h]
\centering
\figplaceholder{14cm}{%
  \textbf{Fig.~A\ref{fig:gallery} — 4-row gallery placeholder.}\\[6pt]
  Row 1 — Static F-formation blocking path.
    Expected: L6 sharp peak over centroid; robot curves around hull.\\[3pt]
  Row 2 — Tight corridor between two dynamic-LF groups.
    Expected: L7 halos nearly touching; planner threads the gap.\\[3pt]
  Row 3 — Mixed crowd (groups + individuals).
    Expected: middle panel shows 3 distinct slot colors for groups, grey
    for individuals.\\[3pt]
  Row 4 — Failure: bilateral flanking.
    Expected: L6+L7 saturate both sides; robot oscillates.\\[6pt]
  \textit{Generate:}
  \texttt{python visualize\_grace.py -{}-model\_dir trained\_models/gram\_map/stageC
    -{}-test\_model 41665.pt -{}-seed \{3,7,11,15\} -{}-max-steps 400}\\
  Extract key frame: \texttt{ffmpeg -i scene.mp4 -vframes 1 -ss \{T\} frame.png}%
}
\caption{\textbf{Cost-map gallery — four representative scenarios.}
  L6 (group cohesion) consistently locates group centroids;
  L7 (repulsion) provides the early-warning margin that routes the planner
  around group hulls without explicit programming.}
\label{fig:gallery}
\end{figure}

%===============================================================================
\section{Slot Assignment Analysis}
\label{sec:slots}

\subsection{Temporal Stability of Slot Assignments}

\begin{figure}[h]
\centering
\figplaceholder{5cm}{%
  \textbf{Fig.~A\ref{fig:slot_stability} — placeholder.}\\[4pt]
  x-axis: timestep (0–200). y-axis: mean assignment entropy
  $H = -\sum_k \alpha_{kn}\log\alpha_{kn}$ averaged over visible humans.\\
  Expected: high entropy ($\approx$0.9) in steps 0–8, stable below 0.3 afterwards.\\[4pt]
  \textit{Generate:} log \texttt{actor\_critic.base.\_last\_alpha} each step
  in \texttt{test.py}; compute entropy per human per step; plot mean $\pm$ std.%
}
\caption{\textbf{Slot assignment entropy over a full episode.}
  SlotAttention converges to stable, low-entropy group assignments within
  5--10 steps and maintains them for the episode, confirming temporal consistency.}
\label{fig:slot_stability}
\end{figure}

\subsection{Learned Assignments vs.\ Ground-Truth Groups}

\begin{figure}[h]
\centering
\figplaceholder{5cm}{%
  \textbf{Fig.~A\ref{fig:slot_gt} — placeholder.}\\[4pt]
  Side-by-side heatmaps at one timestep. Left: $\alpha \in \mathbb{R}^{3\times20}$
  (rows = slots, columns = humans, color = weight). Right: GT binary membership
  matrix (same layout, one bright column-block per group).\\
  Expected: each GT group block aligns with one dominant slot row on the left.\\[4pt]
  \textit{Extract:} \texttt{actor\_critic.base.\_last\_alpha[0].numpy()} in
  \texttt{visualize\_grace.py}; plot with \texttt{plt.imshow}.%
}
\caption{\textbf{Learned slot assignments vs.\ ground-truth group membership.}
  Each group is captured by one dominant slot; individuals receive distributed
  weights, confirming the competitive-softmax behaviour described in Section~3.2.}
\label{fig:slot_gt}
\end{figure}

%===============================================================================
\section{Failure Case Analysis}
\label{sec:failures}

\subsection{Failure Case 1 — Bilateral Group Corridor}

\begin{figure}[h]
\centering
\figplaceholder{5cm}{%
  \textbf{Fig.~A\ref{fig:failure_bilateral} — placeholder.}\\[4pt]
  Three-panel frame from \texttt{visualize\_grace.py} at the collision step.\\
  Scenario: two dynamic-LF groups converge from both sides simultaneously.\\
  Key: L6 (channel 5) shows two peaks bracketing the robot position;
  L7 (channel 6) halos overlap in the robot's forward direction.\\[4pt]
  \textit{Find seed:} sweep \texttt{-{}-seed} 20–40 until a bilateral-flank
  failure is observed in the left panel trajectory.%
}
\caption{\textbf{Failure case 1: bilateral group corridor.}
  Both L6 peaks are active with no navigable corridor between them.
  The robot oscillates and eventually grazes the nearer group hull.}
\label{fig:failure_bilateral}
\end{figure}

\subsection{Failure Case 2 — Dense Crowd Saturation}

\begin{figure}[h]
\centering
\figplaceholder{5cm}{%
  \textbf{Fig.~A\ref{fig:failure_saturation} — placeholder.}\\[4pt]
  Three-panel frame with \texttt{config.sim.human\_num = 28}.\\
  Key: L1 (individual occupancy) near-uniform ($>$0.8); L6 no longer the
  dominant channel — group cohesion signal is buried under individual noise.\\[4pt]
  \textit{Generate:} temporarily set \texttt{human\_num=28} in config,
  run \texttt{test.py}, capture a timeout episode frame.%
}
\caption{\textbf{Failure case 2: cost-map saturation at $N>25$ pedestrians.}
  Individual and trajectory channels saturate, masking the group cohesion layer.
  The planner reverts to individual-occupancy-based steering and times out.}
\label{fig:failure_saturation}
\end{figure}
 in grace.tex)
% Not a standalone document — no \documentclass or \begin{document}.
% ============================================================

% Reset counters so appendix figures/tables are numbered A1, A2, ...
\renewcommand{\thefigure}{A\arabic{figure}}
\renewcommand{\thetable}{A\arabic{table}}
\renewcommand{\theequation}{A\arabic{equation}}
\setcounter{figure}{0}
\setcounter{table}{0}
\setcounter{equation}{0}

% ── Placeholder helper (requires tcolorbox — loaded in grace.tex preamble) ────
\newcommand{\figplaceholder}[2]{%
  \begin{tcolorbox}[colback=gray!10,colframe=gray!50,
                    width=\linewidth,height=#1,
                    halign=center,valign=center]
    \textit{\small #2}
  \end{tcolorbox}%
}

\noindent
This appendix provides full implementation details~(\S\ref{sec:impl}),
group-detection evaluation against DBSCAN baselines~(\S\ref{sec:detection}),
extended ablation analysis with training curves and per-seed
breakdowns~(\S\ref{sec:ablation_ext}), a cost-map visualization
gallery~(\S\ref{sec:gallery}), slot assignment
analysis~(\S\ref{sec:slots}), and failure-case analysis~(\S\ref{sec:failures}).

%===============================================================================
\section{Implementation Details}
\label{sec:impl}

\subsection{Architecture Specifications}

Table~\ref{tab:arch} gives the full specification of every module in \grace{}.

\begin{table}[h]
\centering
\caption{Complete \grace{} architecture with parameter counts.}
\label{tab:arch}
\small
\begin{tabular}{llr}
\toprule
\textbf{Module} & \textbf{Specification} & \textbf{Params} \\
\midrule
\multicolumn{3}{l}{\textit{Perception backbone (pretrained, Stage A)}} \\
GroupDetector   & GNN, 3 layers, input 21-d, embed 64-d & \tbd \\
SlotAttention   & $K=3$ slots, 64-d, 3 iters            & \tbd \\
\midrule
\multicolumn{3}{l}{\textit{Cost-map synthesis (analytic, zero learnable params)}} \\
CostMapSynthesizer & 9 ch, $32{\times}32$ grid, 6\,m range & 0 \\
\midrule
\multicolumn{3}{l}{\textit{Planning head (trained, Stages B--C)}} \\
CostMapPlanner  & Conv 32→64→128→128, AdaptivePool(4)        & \tbd \\
Robot MLP       & Linear 9→128→64                            & \tbd \\
Fusion          & Linear 320→256, ReLU                       & \tbd \\
GRUCell         & 256→256                                    & \tbd \\
Actor head      & Linear 256→256, Tanh $\times$2             & \tbd \\
Critic head     & Linear 256→256, Tanh $\times$2, Linear 256→1 & \tbd \\
OccupancyHead   & Conv $C$→32→$T$ ($1{\times}1$), no sigmoid & \tbd \\
\midrule
\textbf{Total trainable (Stage C)} & & \tbd \\
\bottomrule
\end{tabular}
\end{table}

\subsection{Full Hyperparameter Table}

\begin{table}[h]
\centering
\caption{PPO and curriculum hyperparameters for all training stages.}
\label{tab:hparams}
\small
\begin{tabular}{llll}
\toprule
\textbf{Hyperparameter} & \textbf{Stage A} & \textbf{Stage B} & \textbf{Stage C} \\
\midrule
Algorithm           & Supervised          & PPO              & PPO \\
Learning rate       & $1\!\times\!10^{-3}$ & $5\!\times\!10^{-4}$ & $5\!\times\!10^{-5}$ \\
LR schedule         & Cosine              & Linear decay     & Linear decay \\
Parallel envs       & —                   & 16               & 16 \\
Rollout steps       & —                   & 30               & 30 \\
PPO epochs          & —                   & 5                & 5 \\
Mini-batches        & —                   & 2                & 2 \\
Clip $\epsilon$     & —                   & 0.2              & 0.2 \\
Discount $\gamma$   & —                   & 0.99             & 0.99 \\
GAE $\lambda$       & —                   & 0.95             & 0.95 \\
Entropy coef        & —                   & 0.00             & 0.00 \\
Value loss coef     & —                   & 0.5              & 0.5 \\
Grad clip norm      & —                   & 0.5              & 0.5 \\
Aux loss $\lambda$  & —                   & 0.0 (off)        & 0.1 \\
Backbone frozen     & —                   & Yes              & No  \\
Total env steps     & —                   & 20M              & 20M \\
\midrule
\multicolumn{4}{l}{\textit{Environment — Stage B}} \\
\# humans           & — & 10  & — \\
Circle radius (m)   & — & 5.0 & — \\
Groups              & — & 2×static-F & — \\
On-path groups      & — & 0   & — \\
Realistic phases    & — & Off & — \\
\midrule
\multicolumn{4}{l}{\textit{Environment — Stage C}} \\
\# humans           & — & — & 20 \\
Circle radius (m)   & — & — & 8.5 \\
Groups              & — & — & 3×mixed \\
On-path groups      & — & — & 2 \\
Realistic phases    & — & — & A–E (all on) \\
\bottomrule
\end{tabular}
\end{table}

\subsection{Reward Shaping Details}

The reward at each non-terminal step is $r_t = r_\text{pot} + r_\text{grp} + r_\text{ind}$, where:
\begin{align}
r_\text{pot} &= 2\,(\text{dist}_{t-1} - \text{dist}_t) \label{eq:pot} \\
r_\text{ind} &= (d_\text{thr} - d_\text{ind}) \cdot 10 \cdot \Delta t
  \quad \text{if } d_\text{ind} < 0.25\,\text{m} \label{eq:ind} \\
r_\text{grp} &= (d_\text{grp,thr} - d_\text{grp}) \cdot 10 \cdot \Delta t
  \quad \text{if } d_\text{grp} < 0.35\,\text{m} \label{eq:grp}
\end{align}
$r_\text{pot}$ is always active ($\approx{+}0.5$/step at $v_\text{pref}=1$\,m/s).
Terminal events override $r_t$: success $+10$, collision $-20$, group intrusion $-5$.

%===============================================================================
\section{Group Detection Evaluation}
\label{sec:detection}

To validate the GroupDetector + SlotAttention backbone \emph{independently of navigation},
we evaluate it on 100 held-out episodes from the realistic benchmark.
Ground-truth group membership is read from \texttt{human.group\_id};
predicted pair scores are $\hat{y}_{ij} = \sum_k \alpha_{ki}\alpha_{kj}$.

\begin{table}[h]
\centering
\caption{\textbf{Group detection performance vs.\ DBSCAN baselines.}
DBSCAN uses position-only features; DBSCAN+V uses position and velocity.
Averaged over 100 episodes. \tbd{} = run \texttt{eval\_group\_detection.py} to fill.}
\label{tab:detection}
\small
\begin{tabular}{lccccc}
\toprule
\textbf{Method} & \textbf{Pair F1}$\uparrow$ & \textbf{AUC}$\uparrow$
  & \textbf{Purity}$\uparrow$ & \textbf{ARI}$\uparrow$ & \textbf{NMI}$\uparrow$ \\
\midrule
DBSCAN ($\epsilon$=0.5) & \tbd & \tbd & \tbd & \tbd & \tbd \\
DBSCAN ($\epsilon$=0.8) & \tbd & \tbd & \tbd & \tbd & \tbd \\
DBSCAN ($\epsilon$=1.0) & \tbd & \tbd & \tbd & \tbd & \tbd \\
DBSCAN ($\epsilon$=1.5) & \tbd & \tbd & \tbd & \tbd & \tbd \\
DBSCAN+V ($\epsilon$=0.8) & \tbd & \tbd & \tbd & \tbd & \tbd \\
\midrule
\best{\grace{} backbone (ours)} & \tbd & \tbd & \tbd & \tbd & \tbd \\
\bottomrule
\end{tabular}
\end{table}

\textit{Command:}
\texttt{python eval\_group\_detection.py -{}-model\_dir trained\_models/gram\_map/stageC
-{}-test\_model 41665.pt -{}-episodes 100 -{}-seed 0 -{}-out results/detection\_grace.csv}

\subsection{Detection Robustness Under Occlusion}
\label{sec:detection_occlusion}

We repeat the evaluation dropping 20\%, 40\%, and 60\% of visible humans at random
each step to simulate partial occlusion.

\begin{table}[h]
\centering
\caption{\textbf{Pair F1 under increasing occlusion.}
\grace{} is expected to degrade more gracefully than DBSCAN because the GNN
uses velocity context in addition to position.}
\label{tab:occlusion}
\small
\begin{tabular}{lccc}
\toprule
\textbf{Method} & \textbf{0\% dropped} & \textbf{40\% dropped} & \textbf{60\% dropped} \\
\midrule
DBSCAN (best $\epsilon$)  & \tbd & \tbd & \tbd \\
\grace{} backbone         & \tbd & \tbd & \tbd \\
\bottomrule
\end{tabular}
\end{table}

%===============================================================================
\section{Extended Ablation Analysis}
\label{sec:ablation_ext}

\subsection{Training Curves for All Ablation Variants}

\begin{figure}[h]
\centering
\figplaceholder{7cm}{%
  \textbf{Fig.~A\ref{fig:training_curves} — placeholder.}\\[4pt]
  Single plot: mean PPO reward vs.\ training updates.\\
  Lines: \grace{} full (black), C1 no-group (red), C3 no-traj (blue),
  C5 uniform-$\alpha$ (green), C4 no-aux seeds 0/1/2 (orange dashed).\\[4pt]
  \textit{Generate:} \texttt{python plot.py} — configure \texttt{model\_dirs}
  list to include all ablation output directories.\\
  Key visual: C4 orange lines should diverge widely, all others converge.%
}
\caption{\textbf{Training reward curves for \grace{} and all ablation variants.}
  C4 (no auxiliary loss) shows high variance across seeds; all other ablations
  converge, though at lower final reward than the full model.}
\label{fig:training_curves}
\end{figure}

\subsection{Per-Seed Breakdown for C4 (No Auxiliary Loss)}

\begin{table}[h]
\centering
\caption{\textbf{C4: per-seed navigation results.}
Three Stage~C training runs from the same Stage~B checkpoint, auxiliary loss disabled.
High variance confirms the instability claim in Section~6 of the main paper.}
\label{tab:c4_seeds}
\small
\begin{tabular}{lcccc}
\toprule
\textbf{Run} & \textbf{SR\%}$\uparrow$ & \textbf{CR\%}$\downarrow$
  & \textbf{GCR\%}$\downarrow$ & \textbf{Converged?} \\
\midrule
Seed 0              & \tbd & \tbd & \tbd & \tbd \\
Seed 1              & \tbd & \tbd & \tbd & \tbd \\
Seed 2              & \tbd & \tbd & \tbd & \tbd \\
\midrule
Mean $\pm$ std      & \tbd & \tbd & \tbd & — \\
\midrule
\grace{} full (ref) & $87\pm\tbd$ & $7\pm\tbd$ & $0\pm0$ & Yes \\
\bottomrule
\end{tabular}
\end{table}

\subsection{$K$-Slot Sweep}

\begin{figure}[h]
\centering
\figplaceholder{5cm}{%
  \textbf{Fig.~A\ref{fig:k_sweep} — placeholder.}\\[4pt]
  Dual y-axis line plot. x-axis: $K \in \{1,2,3,4,5\}$.
  Left y-axis: SR\% (inverted-U, peak at $K=3$).
  Right y-axis: GCR\% (U-shape, minimum at $K=3$).\\[4pt]
  \textit{Data:} \texttt{test.py} on
  \texttt{trained\_models/gram\_map/ablation\_K\{1,2,4,5\}} + full model ($K=3$).%
}
\caption{\textbf{Effect of slot count $K$ on SR and GCR.}
  $K=3$ achieves the best trade-off: higher SR than $K\leq2$ and lower GCR than
  $K=1$, with marginal degradation at $K>3$ due to over-segmentation.}
\label{fig:k_sweep}
\end{figure}

%===============================================================================
\section{Cost-Map Visualization Gallery}
\label{sec:gallery}

Fig.~\ref{fig:gallery} shows four critical frames from \texttt{visualize\_grace.py}.
Each row is one scenario: (left) ground-truth environment with group hulls;
(middle) robot's learned slot assignments; (right) all 9 cost-map channels.

\begin{figure}[h]
\centering
\figplaceholder{14cm}{%
  \textbf{Fig.~A\ref{fig:gallery} — 4-row gallery placeholder.}\\[6pt]
  Row 1 — Static F-formation blocking path.
    Expected: L6 sharp peak over centroid; robot curves around hull.\\[3pt]
  Row 2 — Tight corridor between two dynamic-LF groups.
    Expected: L7 halos nearly touching; planner threads the gap.\\[3pt]
  Row 3 — Mixed crowd (groups + individuals).
    Expected: middle panel shows 3 distinct slot colors for groups, grey
    for individuals.\\[3pt]
  Row 4 — Failure: bilateral flanking.
    Expected: L6+L7 saturate both sides; robot oscillates.\\[6pt]
  \textit{Generate:}
  \texttt{python visualize\_grace.py -{}-model\_dir trained\_models/gram\_map/stageC
    -{}-test\_model 41665.pt -{}-seed \{3,7,11,15\} -{}-max-steps 400}\\
  Extract key frame: \texttt{ffmpeg -i scene.mp4 -vframes 1 -ss \{T\} frame.png}%
}
\caption{\textbf{Cost-map gallery — four representative scenarios.}
  L6 (group cohesion) consistently locates group centroids;
  L7 (repulsion) provides the early-warning margin that routes the planner
  around group hulls without explicit programming.}
\label{fig:gallery}
\end{figure}

%===============================================================================
\section{Slot Assignment Analysis}
\label{sec:slots}

\subsection{Temporal Stability of Slot Assignments}

\begin{figure}[h]
\centering
\figplaceholder{5cm}{%
  \textbf{Fig.~A\ref{fig:slot_stability} — placeholder.}\\[4pt]
  x-axis: timestep (0–200). y-axis: mean assignment entropy
  $H = -\sum_k \alpha_{kn}\log\alpha_{kn}$ averaged over visible humans.\\
  Expected: high entropy ($\approx$0.9) in steps 0–8, stable below 0.3 afterwards.\\[4pt]
  \textit{Generate:} log \texttt{actor\_critic.base.\_last\_alpha} each step
  in \texttt{test.py}; compute entropy per human per step; plot mean $\pm$ std.%
}
\caption{\textbf{Slot assignment entropy over a full episode.}
  SlotAttention converges to stable, low-entropy group assignments within
  5--10 steps and maintains them for the episode, confirming temporal consistency.}
\label{fig:slot_stability}
\end{figure}

\subsection{Learned Assignments vs.\ Ground-Truth Groups}

\begin{figure}[h]
\centering
\figplaceholder{5cm}{%
  \textbf{Fig.~A\ref{fig:slot_gt} — placeholder.}\\[4pt]
  Side-by-side heatmaps at one timestep. Left: $\alpha \in \mathbb{R}^{3\times20}$
  (rows = slots, columns = humans, color = weight). Right: GT binary membership
  matrix (same layout, one bright column-block per group).\\
  Expected: each GT group block aligns with one dominant slot row on the left.\\[4pt]
  \textit{Extract:} \texttt{actor\_critic.base.\_last\_alpha[0].numpy()} in
  \texttt{visualize\_grace.py}; plot with \texttt{plt.imshow}.%
}
\caption{\textbf{Learned slot assignments vs.\ ground-truth group membership.}
  Each group is captured by one dominant slot; individuals receive distributed
  weights, confirming the competitive-softmax behaviour described in Section~3.2.}
\label{fig:slot_gt}
\end{figure}

%===============================================================================
\section{Failure Case Analysis}
\label{sec:failures}

\subsection{Failure Case 1 — Bilateral Group Corridor}

\begin{figure}[h]
\centering
\figplaceholder{5cm}{%
  \textbf{Fig.~A\ref{fig:failure_bilateral} — placeholder.}\\[4pt]
  Three-panel frame from \texttt{visualize\_grace.py} at the collision step.\\
  Scenario: two dynamic-LF groups converge from both sides simultaneously.\\
  Key: L6 (channel 5) shows two peaks bracketing the robot position;
  L7 (channel 6) halos overlap in the robot's forward direction.\\[4pt]
  \textit{Find seed:} sweep \texttt{-{}-seed} 20–40 until a bilateral-flank
  failure is observed in the left panel trajectory.%
}
\caption{\textbf{Failure case 1: bilateral group corridor.}
  Both L6 peaks are active with no navigable corridor between them.
  The robot oscillates and eventually grazes the nearer group hull.}
\label{fig:failure_bilateral}
\end{figure}

\subsection{Failure Case 2 — Dense Crowd Saturation}

\begin{figure}[h]
\centering
\figplaceholder{5cm}{%
  \textbf{Fig.~A\ref{fig:failure_saturation} — placeholder.}\\[4pt]
  Three-panel frame with \texttt{config.sim.human\_num = 28}.\\
  Key: L1 (individual occupancy) near-uniform ($>$0.8); L6 no longer the
  dominant channel — group cohesion signal is buried under individual noise.\\[4pt]
  \textit{Generate:} temporarily set \texttt{human\_num=28} in config,
  run \texttt{test.py}, capture a timeout episode frame.%
}
\caption{\textbf{Failure case 2: cost-map saturation at $N>25$ pedestrians.}
  Individual and trajectory channels saturate, masking the group cohesion layer.
  The planner reverts to individual-occupancy-based steering and times out.}
\label{fig:failure_saturation}
\end{figure}
